\naslov{Konvergenca porazdelitev}

\begin{definicija}
  Zaporedje realnih slučajnih spremenljivk $X_1, X_2, \ldots$ \pojem{v
	porazdelitvi konvergira} proti slučajni spremenljivki $X$, če za vsak $x \in
  \R$, za katerega je $P(X = x) = 0$, velja
  \[
	\lim_{n \to \infty} P(X_n \le x) = P(X \le x).
  \]
  Pišemo $X_n \xrightarrow[n \to \infty]{d} X$.
\end{definicija}

\begin{opomba}
  Temu pravimo tudi \pojem{šibka konvergenca}.
\end{opomba}

\begin{definicija}
  Zaporedje slučajnih spremenljivk $X_1, X_2, \ldots$ z vrednostmi v topološkem
  prostoru $S$ \pojem{v porazdelitvi konvergira proti $X$}, če za vsako zvezno
  in omejeno funkcijo $h: S \to \R$ velja
  \[
	\lim_{n \to \infty} E(h(X_n)) = E(h(X)).
  \]
  Funkcijam $h$ rečemo tudi \pojem{preizkusne} ali \pojem{testne} funkcije.
\end{definicija}

\begin{izrek}[Helly-Bray]
  Za porazdelitve na realni osi z običajno topologijo definiciji sovpadata.
\end{izrek}

\begin{izrek}[Centralni limitni izrek]
  Če so $X_1, X_2, \ldots$ neodvisne in enako porazdeljene z $E(X_1) = \mu_1$ in
  $\var(X_1) = \sigma_1^2$, velja
  \[
	\frac{X_1 + \ldots + X_n - n \mu_1}{\sqrt{n}} \xrightarrow[n\to \infty]{d}
	N(0, \sigma_1^2).
  \]
\end{izrek}

\vprasanje{Definiraj konvergenco v porazdelitvi in reformuliraj centralni
  limitni izrek.}

\begin{trditev}
  Naj bodo $X_1, X_2, \ldots, X$ slučajne spremenljivke z vrednostmi v
  topološkem prostoru $S$ in $T$ še en topološki prostor.
  Naj bo $g: S \to T$ zvezna.
  Če velja $X_n \xrightarrow[n \to \infty]{d} X$, velja tudi $g(X_n)
  \xrightarrow[n\to\infty]{d} g(X)$.
\end{trditev}

\begin{proof}
  Naj bo $h : T \to \R$ zvezna in omejena.
  Tedaj je tudi $h \circ g : S \to \R$ zvezna in omejena, zato je
  \[
	E(h(g(X))) = \lim_{n \to \infty} E(h(g(X_n))) = \lim_{n \to \infty} E(h
	\circ g(X_n)) = E(h \circ g(X)).
  \]
\end{proof}

\begin{trditev}
  Naj bo $S_0 \subseteq S$ in naj imajo $X_1, X_2, \ldots, X$ vrednosti v $S_0$.
  Če gre $X_n \xrightarrow[n \to \infty]{d} X$ v okviru $S_0$, to velja tudi v
  okviru prostora $S$.
\end{trditev}

\begin{proof}
  Naj bo $h: S \to \R$ zvezna in omejena, ter naj bo $h_0$ zožitev $h$ na $S_0$.
  Tudi ta funkcija je zvezna in omejena, zato je $\lim_{n \to \infty}
  E(h_0(X_n)) = E(h_0(X))$.
\end{proof}

\begin{trditev}
  Naj bo $S$ metrizabilen prostor, $S_0 \subseteq S$ pa zaprt podprostor.
  Če imajo $X_1, X_2, \ldots, X$ vrednosti v $S_0$ in $X_n \xrightarrow[n \to
  \infty]{d} X$ v okviru $S$, to velja tudi v okviru $S_0$.
\end{trditev}

\begin{proof}
  Vsaka zvezna omejena funkcija $h_0 : S_0 \to \R$ se da po Tietzejevem izreku
  razširiti do zvezne in omejene funkcije $h: S \to \R$.
\end{proof}

\vprasanje{V katerem primeru se konvergenca v porazdelitvi razširi iz večjega
  prostora v podprostor? Kaj je ideja dokaza?}

\begin{izrek}
  Naj bo $S$ metrizabilen in $S_0 \subseteq S$ odprt podprostor.
  Naj bodo $X_1, X_2, \ldots$ slučajne spremenljivke z vrednostmi v $S$ in $X$ z
  vrednostmi v $S_0$. Naj gre $X_n \xrightarrow[n\to\infty]{d} X$ v okviru
  prostora $S$.
  Nadalje naj bodo $X_1^*, X_2^*, \ldots$ slučajne spremenljivke z vrednostmi v
  $S_0$ ter naj bo $X_n^*(\omega) = X_n(\omega)$ brž ko je $X_n(\omega) \in
  S_0$.
  Tedaj gre $X_n \xrightarrow[n\to\infty]{d} X$ tudi v okviru $S_0$.
\end{izrek}

\begin{trditev}
  Naj bodo $X_1, X_2, \ldots$ slučajne spremenljivke z vrednostmi v metričnem
  prostoru $(S, d)$.
  Naj bo $c \in S$ konstantna.
  Tedaj gre $X_n \xrightarrow[n \to \infty]{d} c$ natanko tedaj, ko za vsak $r >
  0$ velja
  \[
	\lim_{n \to \infty} P(d(X_n,c) \ge r) = 0.
  \]
\end{trditev}

\begin{proof}
  V desno:
  Obstaja taka zvezna funkcija $h: S \to [0,1]$, da je $h(c) = 0$ in $h(x) = 1$
  za $d(x, c) \ge r$.
  Velja $\mathbbm{1}(d(x, c) \ge r) \le h(x)$, na tej neenakosti pa lahko
  uporabimo pričakovano vrednost in dobimo
  \[
	P(d(X_n, c) \ge r) \le E(h(X_n)) \xrightarrow[n \to \infty]{} E(h(c)) = 0.
  \]

  V levo:
  Naj bo $h: S \to [0,1]$ zvezna in $\varepsilon > 0$.
  Obstaja tak $r > 0$, da za vsak $x \in S$ z $d(x, c) \ge r$ velja $\abs{h(x) -
	h(c)} < \nicefrac{\varepsilon}{2}$.
  Sledi
  \begin{align*}
	&\abs{E(h(X_n)) - h(c)} \\
	&\le E(\abs{h(X_n) - h(c)}) \\
	&= E(\abs{h(X_n) - h(c)} \mathbbm{1}(d(X_n, c) < r))
	  + E(\abs{h(X_n) - h(c)} \mathbbm{1}( d(X_n, c) \ge r )) \\
	&\le \frac{\varepsilon}{2} + P(d(X_n, c) \ge r).
  \end{align*}
  Za dovolj pozne $n$ bo to manj od $\varepsilon$.
\end{proof}

\vprasanje{Kdaj zaporedje slučajnih spremenljivk konvergira k konstanti?
  Dokaži.}

\begin{trditev}[Neenačba Markova]
  Za nenegativno slučajno spremenljivko $W$ in poljuben $a > 0$ velja
  \[
	P(W \ge a) \le \frac{E(W)}{a}.
  \]
\end{trditev}

\begin{proof}
  Velja $\mathbbm{1}(w \ge a) \le \nicefrac{w}{a}$.
  Na tem uporabimo pričakovano vrednost.
\end{proof}

\vprasanje{Povej in dokaži neenačbo Markova.}

\begin{trditev}
  Če je $p > 0$ in $\lim_{n \to \infty} E({(d(X_n, c))}^p) = 0$, gre $X_n
  \xrightarrow[n \to \infty]{d} c$.
\end{trditev}

\begin{proof}
  Velja $P(d(X_n, c) \ge r) = P({(d(X_n, c))}^p \ge r^p)$.
  Na drugem delu uporabimo neenačbo Markova.
\end{proof}

\begin{trditev}[šibki zakon velikih števil]
  Če so $X_1, X_2, \ldots$ neodvisne in enako porazdeljene z $E(X_1^2) < \infty$
  in $E(X_1) = \mu$, gre
  \[
	\frac{X_1 + \cdots + X_n}{n} \xrightarrow[n \to \infty]{d} \mu.
  \]
\end{trditev}

\begin{proof}
  Pričakovana vrednost izraza je očitno $\mu$.
  Potem je
  \[
	E\!\left({\left( \frac{X_1 + \cdots + X_n}{n} - \mu \right)}^2\right)
	= \var\left(\frac{X_1 + \cdots + X_n}{n}\right)
	= \frac{\var(X_1)}{n},
  \]
  to pa konvergira k $0$ za $n \to \infty$.
\end{proof}

\vprasanje{Povej in dokaži šibki zakon velikih števil.}

\begin{opomba}
  Velja celo
  \[
	P\!\left( \lim_{n \to \infty} \frac{X_1 + \cdots + X_n}{n} = \mu \right) = 1.
  \]
  Temu pravimo \pojem{krepki zakon velikih števil}.
\end{opomba}

\begin{izrek}
  Naj bo $S$ metrizabilen prostor, ki je števna unija svojih kompaktnih
  podprostorov.
  Naj bodo $X_1, X_2, \ldots$ slučajne spremenljivke z vrednostmi v $S$.
  Naj bo $T$ še en metrizabilen prostor in naj imajo $Y_1, Y_2, \ldots$
  vrednosti v $T$.
  Če gre $X_n \xrightarrow[n \to \infty]{d} X$ in $Y_n \xrightarrow[n \to
  \infty]{d} c$, kjer je $c \in T$ konstanta, gre tudi $(X_n, Y_n)
  \xrightarrow[n \to \infty]{d} (X, c)$.
\end{izrek}

\vprasanje{Kaj velja za konvergenco parov slučajnih spremenljivk?}